\documentclass[a4paper,14pt]{exam}
\usepackage{array,amsmath,fancybox,amsfonts,bm}
\usepackage{graphicx,amssymb,wasysym,latexsym,slashed,anysize,subfigure,calc}



%versores: %uso: \uveci\ne\hat{\imath}_{\uvec{i}}+\uvec{j}+\uvec{k}$
	\usepackage{bm}
	\newcommand{\uveci}{{\bm{\hat{\textnormal{\bfseries\i}}}}}
	\newcommand{\uvecj}{{\bm{\hat{\textnormal{\bfseries\j}}}}}
	\DeclareRobustCommand{\uvec}[1]{{%
 	 \ifcsname uvec#1\endcsname
  	   \csname uvec#1\endcsname
  	 \else
   		 \bm{\hat{\mathbf{#1}}}%
   		\fi
									}}



%tamaños fuentes extra: 8,9,14,17,20
	\usepackage{extsizes}
	
 %caja bonita
  \usepackage[many]{tcolorbox}
\definecolor{titlebg}{RGB}{80,100,72}
\newtcolorbox{learn}{
  breakable,
  enhanced,
  arc=0pt,
  outer arc=0pt,
  colframe=titlebg,
  colback=titlebg!7,
  overlay unbroken and first={
    \node[
      draw=titlebg,
      fill=titlebg,
      rotate=270,
      anchor=north west,
      text=white,
      font=\bfseries
    ]
    at (frame.north west)  
    {DATOS};
  }
}


\newtcolorbox{problemas}{
  breakable,
  enhanced,
  arc=0pt,
  outer arc=0pt,
  colframe=titlebg,
  colback=titlebg!7,
  overlay unbroken and first={
    \node[
      draw=titlebg,
      fill=titlebg,
      rotate=270,
      anchor=north west,
      text=white,
      font=
    ]
    at (frame.north west)  
    {AVISO};
  }
}







%imágenes fundidas texto
\usepackage{wrapfig}

\usepackage{mwe,float,afterpage}

\usepackage{float}
	%\marginsize{left}{right}{top}{bottom}
    \marginsize{.6cm}{1cm}{1.2cm}{.1cm}
    
    
    
    
%entorno para química: texto y math mode
%\newcommand*\chem[1]{\textup{#1}}
\newcommand*\chem[1]{\ensuremath{\mathrm{#1}}}

%símbolo bonito para qed
\usepackage{amsthm}
	%\newcommand{\mc}[3]{\multicolumn{#1}{#2}{#3}}
\newcommand{\halmos}{
					\begin{flushright}
						\rule{1mm}{2.5mm} 
					\end{flushright}
					  }
\renewcommand{\qed}{\halmos}



\usepackage[spanish]{babel}
\usepackage[utf8]{inputenc}
%\usepackage{fancyhdr}
\usepackage{lettrine}

%subrayado y highlight. Tiene que estar detrás de usepackage[color]
\usepackage{soul}
	%\textst{normal} 
	\DeclareRobustCommand{\hlcyan}[1]{{\sethlcolor{cyan}\hl{#1}}} %en cyan!
		%\hl{highlighted text} %uso


%fuentes
	\usepackage[T1]{fontenc}
	\usepackage{mathpazo}

	\usepackage{librebaskerville}

%fancy tables
	\usepackage{booktabs}
	\usepackage[small,centerlast,up]{caption}

	

\usepackage{amssymb,amsmath}
\newcommand{\angstrom}{\mbox{\normalfont\AA}}

\usepackage{graphicx,graphics}
\usepackage{color}
	%	\colorgrids %grids con color
	%	\definecolor{GridColor}{gray}{0.8}
		%grid
%			\setlength{\gridsize}{5mm}
%   			 \setlength{\gridlinewidth}{0.1pt}

%fracciones en texto más elaboradas
	\usepackage{xfrac}

%mejora en la exposición de guarismos
	\usepackage{numprint}
 


%redefiniciones
\pointpoints{punto}{puntos}
\bonuspointpoints{punto}{puntos}
\hqword{Pregunta} %substitutes text for \Question"
%\vpgword{Pág.}%substitutestext for \Page"
\hpword{Máx.}%substitutes text for \Points"
\hsword{Nota}%substitutes text for \Score"
\htword{Total}%substitutes text for \Total:"

\chqword{Pregunta}
\chpword{Máx.}
\chsword{Nota}
\chtword{Total}
\chbpword{Puntos \textit{bonus}}

\renewcommand{\solutiontitle}{\noindent\textbf{Solución:}\par\noindent}




%para los entornos de soluciones
\usepackage{mdframed,xcolor,color}
	%uso: 	entorno (begin,end): {mdframed}[backgroundcolor=brown!10] 


\usepackage{everypage,draftwatermark} %Marca de agua
\SetWatermarkText{\parbox{2\textwidth}{\centering B2\\dist.}}
\SetWatermarkScale{1.5}
\SetWatermarkLightness{0.92}





\pagestyle{headandfoot}
%\runningheadrule
%\firstpageheader{Física para Biólogos}{Boletín 1a}{1 de octubre de 2011}
%\runningheader{Física para Biólogos}
%              {Boletín 1a, Página \thepage\ de \numpages}
%            {1 de octubre de 2011}
%\firstpagefooter{}{}{}
%\runningfooter{}{}{}


\cfoot{}
%espacios extra
%\extrawidth{1in}
%%\extrafootheight{.75in}
%\extraheadheight{1in}

%cabecera
\header{\includegraphics[scale=.22]{a.png}}       
      {}
       {$f_2$\textbf{B2-{\small distancia}}\\\textbf{febrero}{\small 2024}}       
%pie de página

   %enumera en todas menos en la primera
   %\runningfootrule

   %enumera en todas
   %\footrule
   %\footer{}{Página \thepage\ de \numpages}{}
   \footer{}
	{%Página \thepage\ de \numpages
	}
	{%\iflastpage{{\footnotesize \textit{Fin del examen}}}{{\footnotesize \textit{Da la vuelta a la hoja}\ldots}}
	}

	
\begin{document}
	
\fbox{\parbox{0.945\linewidth}{\textbf{Nombre y apellidos}:\vspace*{20mm}}}

\vspace*{1cm}

\begin{minipage}{.07\linewidth}

\end{minipage}\hfill
\begin{minipage}{.93\linewidth}
\begin{problemas}{Un examen es una \textit{comunicación}: debes invocar las leyes en que te basas para resolver los problemas e \textit{indicar brevemente los razonamientos} que sigues para obtener los resultados.}
\end{problemas}
\end{minipage}

\vspace{1ex}

%\hspace*{1cm}\begin{learn}
%{$\triangleright\; M(\chem{O})\simeq16\;g\,mol^{-1}$\newline
%$\triangleright\; M(\chem{N})\simeq14\;g\,mol^{-1},$\newline
%$\triangleright\; M(\chem{H})\simeq1\;g\,mol^{-1},$\newline
%$\triangleright\; M(\chem{S})\simeq32\;g\,mol^{-1},$\newline
%$\triangleright\; M(\chem{Cu})\simeq64\;g\,mol^{-1},$\newline
%$\triangleright\; R\simeq 0.082\;atm\,L\,mol^{-1}\,K^{-1}$}
%\end{learn}	



\newpage
%\qformat{\textbf{Pregunta \thequestion}\quad (\thepoints)\hfill} 
	
\begin{questions}

	\begin{question}
	\begin{parts}
			\part Describe en qué consiste la \textbf{hipermetropía} y \textbf{cómo se corrige}.
			\begin{mdframed}[backgroundcolor=yellow!20]  
			La hipermetropía es un defecto visual que impide ver bien objetos cercanos. Se produce cuando los rayos de luz convergen por detrás de la retina y se corrige con lentes convergentes.
			\end{mdframed}
			
			\part Una persona \textit{miope} utiliza gafas con lentes divergentes de $5$ dioptrías  para ver correctamente objetos muy alejados.  Estima la \textbf{distancia máxima} a la que puede enfocar nítidamente sin gafas.
			\begin{mdframed}[backgroundcolor=yellow!20]  
			Si las gafas que usa son las correctas, son capaces de traer cualquier objeto (idealmente $s\to-\infty$ a una distancia más cercana que el punto remoto. La ecuación de las lentes delgadas nos dice que debe cumplirse
			\[\frac{1}{f'}=\frac{1}{PR}-\lim_{s\to -\infty}\frac{1}{s}\rightarrow -5=\frac{1}{PR},\]
			de donde $PR\simeq -0.2\;m=-20\;cm.$ 
			
			\end{mdframed}
		\end{parts}				
		\qed
		\end{question}
		
		
		\newpage
		
		\begin{question}
		Un asteroide efectúa orbita circularmente alrededor de una estrella aproximadamente esférica con una energía total de $-10^{10}\; J$. Determine:
		\begin{parts}
			\part El \textbf{cociente entre las energías potencial y cinética} del asteroide.
			\begin{mdframed}[backgroundcolor=yellow!20]  
			La energía cinética del asteroide es
			\(E_c=\sfrac{1}{2}mv^2,\) pero en su movimiento alrededor de la estrella la fuerza centrípeta es la fuerza gravitatoria, así que ha de cumplirse
			\[G\frac{Mm}{r^2}=\frac{mv^2}{r},\]
			que permite escribir $mv^2=\dfrac{GMm}{r}$.	
			Volviendo a la energía cinética resulta entonces
			\[E_c=\frac{1}{2}mv^2=\frac{1}{2}\frac{GMm}{r}=-\frac{1}{2}E_p.\]
			O sea que finalmente la relación entre ambas es $$\frac{E_p}{E_c}=-2.$$
			\end{mdframed}
			\part Los \textbf{valores de ambas energías}: potencial y cinética.
			\begin{mdframed}[backgroundcolor=yellow!20]  
			Sabemos que $E_c+E_p=E_T$ y que $2E_c+E_p=0$. Resolviendo el sistema es posible deducir que
			$$E_c=-E_T=10^{10}\;J,\hspace{4ex}E_p=-2E_c=2E_T=2\cdot 10^{10}\;J.$$
			\end{mdframed}
		\end{parts}
	\qed
		\end{question}
		
	\newpage
	
	
		\begin{question}
		\begin{parts}
			\part La Tierra orbita en torno al Sol describiendo una órbita aproximadamente circular a una distancia de $1\,\text{\footnotesize UA}$. Estima la \textbf{masa del Sol.}
			\begin{mdframed}[backgroundcolor=yellow!20]  
			Utilizando la ley de Kepler podemos ver que 
			\[\frac{R^3}{T^2}=\frac{GM}{4\pi^2}\rightarrow M=\frac{4\pi^2R^3}{GT^2}=4\pi^2\frac{(1.5\cdot 10^{11})^3}{6.67\cdot 10^{-11}\cdot (365\cdot 24\cdot 3600)^2}\simeq 2\cdot 10^{30}\;kg.\]
			\end{mdframed}
			\part Júpiter es un planeta gigante gaseoso situado unas $5.2$ veces más lejos del Sol que la Tierra, ¿\textbf{cuántos años terrestres dura el año joviano}?
		\end{parts}	
	{
	\small \textbf{$\triangleright\;$Datos:} \hspace*{.5cm} $\triangleright\; 1\;\text{\footnotesize año}\simeq 365.25\;\text{\footnotesize días}\hspace{.3cm}
\triangleright\;1\;\text{\scriptsize UA}\simeq 1.5\cdot 10^8\;km\hspace{.3cm} \triangleright\;G\!\simeq\! 6.67\cdot 10^{-11}\;N\,m^2\,kg^{-2}$
	}	
	
	\begin{mdframed}[backgroundcolor=yellow!20]  
	Es muy interesante utilizar un sistema de unidades en el que el tiempo se mida en años y la distancia en unidades astronómicas. En él la ley de Kepler se escribe \(\omega^2R^3=1,\) y entonces
	para los dos planetas del Sistema Solar se cumplirá
	\[\omega^2_TR^3_T=\omega^2_JR^3_J\longrightarrow \frac{T_J}{T_T}=\left(\frac{R_J}{R_T}\right)^{\sfrac{3}{2}}=5.2^{\sfrac{3}{2}}\simeq 11.9.\]
	O sea: el año joviano equivalen a aproximadamente $12$ años terrestres.
	\end{mdframed}
	\end{question}
	\qed
	
	\newpage
	
	
	\begin{minipage}{.61\linewidth}
\begin{question}
 Tres cargas eléctricas estáticas se sitúan en el plano $xy$ como en la figura.  
\begin{itemize}
			\item[(a)] Calcula el \textbf{potencial eléctrico total} en el punto $P(L,L)$.
			\item[(b)] ¿Cuál sería el \textbf{trabajo} necesario \textit{para cambiar de posición} a la carga $\hat{q}$ desde donde está hasta el punto $P(L,L)$?

%{\small \fbox{\textbf{Datos:} $\triangleright\;L=1.2\;cm\hspace{.5cm}\triangleright\;q_1=q_2=5\;nC,\hspace{.7cm} \triangleright\;q_3=-5\; nC$.
%			}}
			
			\end{itemize}
\end{question}	
{\small \textbf{$\triangleright\;$Indicación:} Deja la solución en términos de $L$, $q$ y $\hat q$.}
		
\end{minipage}\hfill
\begin{minipage}{.4\linewidth}

	\hspace*{.5cm}\def\svgwidth{0.9\linewidth}
	\input{triang.pdf_tex}
\end{minipage}
\begin{mdframed}[backgroundcolor=yellow!20] 
\begin{itemize}
	\item[a)] El potencial en $P(L,L)$ se calcula superponiendo las tres contribuciones procedentes de cada carga puntual. Con todo
			\[V\bigg|_{P(L,L)}=K\frac{|q|}{L}+K\frac{|q|}{\sqrt{2}L}-K\frac{|\hat{q}|}{L}\rightarrow \boxed{V\bigg|_{P(L,L)}=\frac{K}{L}\left[|q|\left(1+\frac{1}{\sqrt{2}}\right)-|\hat{q}|\right].}\]
			
	\item[b)] Al cambiar de posición la carga $\hat{q}$ encontramos \textit{la misma configuración} de cargas que al inicio. De esta invariancia se deduce que el proceso inverso\footnote{Cambiar de posición la carga $\hat{q}$ desde $P(L,L)$ hasta el punto que ocupa en la figura.} debe llevar asociado un trabajo idéntico, de donde
	$$W_{i\to f}=W_{f\to i}.$$
	
	Pero también sabemos que las funciones de estado asociadas a un proceso y a su inverso son iguales salvo un signo, de modo que igualmente debe satisfacerse $W_{i\to f}=-W_{f\to i}$. Ambas ecuaciones dan lugar a una única solución trivial $\boxed{W_{i\to f}=W_{f\to i}=0.}$
	
	Se podría haber razonado a partir de que $W=-\Delta E_p$. El resultado habría sido el mismo puesto que al resultar la misma configuración de cargas, la energía potencial al inicio y al final es idéntica.
\end{itemize} 
			
			\end{mdframed}

	\qed


\newpage

\begin{question}
		
		Una fuente luminosa emite luz de longitud $550\;nm$. Se hace incidir la luz sobre un metal produciendo la emisión de electrones. La energía mínima de extracción de los electrones de la red metálica es de $2.10\;eV$. 

\begin{parts}
			\part Calcula la \textbf{velocidad} de los electrones emitidos.
			\begin{mdframed}[backgroundcolor=yellow!20]  
			El balance energético del efecto fotoeléctrico permite escribir
			\[hf=W+\frac{1}{2}mv^2\longrightarrow v=\sqrt{\frac{2}{m}\left(h\frac{c}{\lambda}-W\right)}\simeq 2\cdot 10^5\;ms^{-1}.\] 
			\end{mdframed}
			\part  Calcula el \textbf{potencial} necesario para anular la corriente de electrones.
			\begin{mdframed}[backgroundcolor=yellow!20]  
			El voltaje de parada, $V_s$, debe ser tal que contrarreste la energía cinética de la corriente, \textit{i.e.}
			\[eV_s=\frac{1}{2}mv^2\longrightarrow V_s\simeq 0.11\,V,\]
			\end{mdframed}
		\end{parts}
		\vspace{1ex}
		
	{\small \textbf{$\triangleright\;$Datos:}
	\hspace*{.5cm}$\triangleright\;c_0\simeq 3\cdot 10^8\;ms ^{-1}\hspace{.2cm}\triangleright\;m_e\simeq 9.1\cdot 10^{-31}\;kg\hspace{.2cm}\triangleright\;1\;e\simeq 1.6\cdot 10^{-19}\;C\hspace{.2cm}\triangleright\;h\simeq 6.626\cdot 10^{-34}\;J\,s\hspace{.2cm}
	%\triangleright\;\hbar c\simeq 197\;eV\,nm
	$
	}
	\qed
	\end{question}


 \end{questions}
\end{document}

